\typeout{NT FILE chapter1.tex}%

\chapter{Introduction}
\label{cha:introduction}

\section{Context}

Theory of Computation is a fundamental area in computer science that includes various
fields of study, including this work's focus: Formal Languages and Automata
Theory (FLAT). These play a critical role in modeling computational processes, verifying system behavior, 
and designing compilers and interpreters. 

Despite its theoretical significance, FLAT concepts can be difficult to grasp due to their abstract nature. For students, interactive
visualization and step-by-step simulation are critical aids. OCamlFLAT/OFLAT, developed at NOVA School of Science and Technology, already
supports finite automata, grammars, pushdown automata, and Turing machines. However, it lacks finite-state transducers (FST), a
model that adds output generation to automata and is central to translation, lexical analysis, and many teaching examples.

While FSTs are well-established in the literature, pedagogical tools that let learners build, simulate, and analyze them with
output tracking are scarce or incomplete. This dissertation fills that gap by extending both the OCamlFLAT library and the OFLAT
web interface with native FST support, aiming to deliver a consistent, interactive learning experience alongside the existing
models.

\section{Objectives}

The primary goal of this thesis is to design and implement support for finite-state transducers within the OCamlFLAT/OFLAT ecosystem. This includes:
\begin{itemize}
  \item Defining data structures and algorithms to model FSTs (including acceptance with expected outputs, determinism checks, and minimization).
  \item Exposing output-aware operations in OCamlFLAT (e.g., \texttt{acceptOut}, \texttt{acceptExpectedOutput}, \texttt{inputEdges}) 
  consistent with existing automata APIs.
  \item Integrating FST editing and step-by-step execution into OFLAT with input:output edge labels, accumulated-output display, 
  and nondeterministic branch visuals.
  \item Enabling exercise checking with \texttt{input->output} examples so UI feedback matches library semantics.
\end{itemize}

By enriching the system with transducer functionality, we aim to provide a more comprehensive pedagogical platform that supports a broader 
spectrum of the Chomsky hierarchy and theoretical computation models. This will enhance both teaching and learning by enabling students 
and instructors to explore complex language transformations through hands-on experimentation.

\section{Contributions}
\label{sec:intro_contrib}

The main contributions of this dissertation are:
\begin{itemize}
  \item A finite-state transducer extension to OCamlFLAT with output-aware acceptance, determinism checks, and minimization utilities reused by the UI.
  \item An OFLAT integration that supports interactive FST construction, step-by-step execution with accumulated outputs, branch highlighting, 
  popper counts, and exercise checking for \texttt{input->output} cases.
  \item An evaluation covering unit tests, integration checks in the browser, and representative FST use cases, with guidance for future improvements.
\end{itemize}

\section{Structure of the Document}

The remainder of this document is structured as follows. 
\begin{itemize}
  \item Chapter~\ref{cha:background} introduces the necessary concepts in formal languages and automata, focusing on finite-state transducers.
  \item Chapter~\ref{cha:related_work} surveys related tools and positions this work against existing automata/FST teaching environments.
  \item Chapter~\ref{cha:current_system} presents the OCamlFLAT/OFLAT platform as it existed before this work.
  \item Chapter~\ref{cha:technologies} outlines the key technologies used (OCaml, \texttt{js\_of\_ocaml}, Cytoscape.js, and supporting libraries).
  \item Chapter~\ref{cha:transducer-implementation} details the OCamlFLAT FST implementation (data types, acceptance, determinism, minimization).
  \item Chapter~\ref{cha:oflat-design} describes the OFLAT design and user experience for FSTs (interaction model, visuals, feedback surfaces).
  \item Chapter~\ref{cha:oflat-implementation} explains the OFLAT implementation, including event wiring, data flow, rendering, and exercise checking.
  \item Chapter~\ref{cha:evaluation} presents the evaluation methodology, test layers, and representative use cases.
  \item Chapter~\ref{cha:conclusion} summarizes the results and outlines future work.
\end{itemize}

